\subsection{Helm}

Để hiện thực hóa mô hình Quản lý gói cho Florus, đồ án lựa chọn sử dụng Helm – công cụ được xem là tiêu chuẩn công nghiệp trong việc đóng gói và triển khai ứng dụng trên nền tảng Kubernetes.

\subsubsection{Cấu trúc và nguyên lý hoạt động} Sức mạnh cốt lõi của Helm nằm ở cấu trúc đóng gói được gọi là Helm Charts. Một Chart là tập hợp các tệp cấu hình Kubernetes được tổ chức theo cấu trúc thư mục chuẩn, bao gồm: 
\begin{itemize} 
    \item \textbf{Templates:} Chứa các tệp mô tả tài nguyên (Deployment, Service, Ingress...) dưới dạng khuôn mẫu, sử dụng ngôn ngữ Go Template. 
    \item \textbf{values.yaml:} Chứa các giá trị cấu hình mặc định, tách biệt dữ liệu khỏi logic, cho phép tùy biến tham số mà không can thiệp vào mã nguồn. 
    \item \textbf{Chart.yaml:} Tệp siêu dữ liệu định danh tên, phiên bản và các thông tin bảo trì của gói. 
\end{itemize}

\subsubsection{Chiến lược tổ chức hệ thống: Mô hình Umbrella Chart} 
Điểm nhấn trong việc áp dụng Helm cho hệ thống Florus là việc sử dụng mô hình thiết kế Umbrella Chart. Đối mặt với độ phức tạp của một hệ thống phân tán, việc quản lý hàng chục chart riêng lẻ sẽ gây khó khăn trong việc đồng bộ. Umbrella Chart giải quyết vấn đề này bằng cách hoạt động như một "chart cha", không chứa các template tài nguyên trực tiếp mà định nghĩa toàn bộ hệ thống thông qua danh sách các tài nguyên phụ thuộc.

Dựa trên mô hình này, kiến trúc triển khai của Florus được phân tách thành 3 phân vùng logic rõ ràng nằm dưới sự quản lý của một Umbrella Chart duy nhất: 
\begin{itemize} 
    \item \textbf{Tầng ứng dụng:} Chứa các sub-charts cho các microservices nghiệp vụ cốt lõi của Florus (Api Gateway, Auth Service,...). Việc gom nhóm này giúp quy trình cập nhật logic nghiệp vụ được thực hiện độc lập và nhanh chóng. 
    \item \textbf{Tầng hạ tầng:} Quản lý các dịch vụ stateful và tài nguyên nền tảng cần thiết như Cơ sở dữ liệu (MinIO, MongoDB) và Message Queue (Kafka). Việc tách biệt tầng này đảm bảo sự ổn định của dữ liệu, tránh ảnh hưởng khi cập nhật tầng ứng dụng.
    \item \textbf{Tầng giám sát:} Tập trung các công cụ giảm sát như Prometheus, Grafana để thu thập metrics và logs từ hai tầng trên.
\end{itemize}

\subsubsection{Lợi ích triển khai thực tế} Việc kết hợp Helm và mô hình Umbrella Chart mang lại hiệu quả vận hành cụ thể: 
\begin{itemize} 
    \item \textbf{Triển khai đồng bộ:} Cho phép khởi tạo toàn bộ hệ sinh thái Florus phức tạp (bao gồm cả App, Infra và Monitoring) chỉ với một câu lệnh duy nhất, đảm bảo tính toàn vẹn của các mối liên kết giữa các dịch vụ.
    \item \textbf{Quản lý môi trường linh hoạt:} Sự khác biệt cấu hình giữa môi trường Development và Production cho cả 3 tầng được kiểm soát hoàn toàn thông qua các tệp values riêng biệt. 
\end{itemize}